% ============================================================================
%  3.5 IMPLEMENTACIÓN DEL BACKEND (Hooks de React Query)
% ============================================================================

\section{Implementaci\'on del backend}

\subsection{Hooks personalizados}

Se cre\'o el archivo \texttt{src/hooks/useEventRegistration.ts} con
cinco hooks que encapsulan toda la l\'ogica de comunicaci\'on con
Supabase:

\begin{enumerate}
  \item \textbf{useIsRegistered(userId, eventId):} Verifica si el
    usuario est\'a inscrito en un evento espec\'ifico.
  \item \textbf{useRegisterForEvent():} Inscribir al usuario en un
    evento.
  \item \textbf{useCancelRegistration():} Cancelar la inscripci\'on
    del usuario.
  \item \textbf{useMyRegistrations(userId):} Obtener todas las
    inscripciones activas del usuario.
  \item \textbf{useRegistrationCounts(eventIds):} Obtener el conteo
    de inscripciones para una lista de eventos.
\end{enumerate}

\subsection{Funciones de API}

Cada hook internamente llama a una funci\'on de API que interact\'ua
directamente con Supabase:

\subsubsection*{Verificar inscripci\'on}

\begin{verbatim}
async function checkRegistration(
  userId: string, eventId: string
): Promise<EventRegistration | null> {
  const { data, error } = await supabase
    .from('event_registrations')
    .select('*')
    .eq('user_id', userId)
    .eq('event_id', eventId)
    .in('status', ['confirmed', 'waitlist'])
    .maybeSingle();

  if (error) throw new Error(error.message);
  return data as EventRegistration | null;
}
\end{verbatim}

\subsubsection*{Inscribir en evento}

\begin{verbatim}
async function registerForEvent(
  userId: string, eventId: string
): Promise<EventRegistration> {
  // Verificar si ya está inscrito
  const existing = await checkRegistration(userId, eventId);
  if (existing) {
    throw new Error('Ya estás inscrito en este evento');
  }

  const { data, error } = await supabase
    .from('event_registrations')
    .insert({
      user_id: userId,
      event_id: eventId,
      status: 'confirmed',
    })
    .select()
    .single();

  if (error) throw new Error(error.message);
  return data as EventRegistration;
}
\end{verbatim}

\subsubsection*{Cancelar inscripci\'on}

\begin{verbatim}
async function cancelRegistration(
  userId: string, eventId: string
): Promise<void> {
  const { error } = await supabase
    .from('event_registrations')
    .delete()
    .eq('user_id', userId)
    .eq('event_id', eventId);

  if (error) throw new Error(error.message);
}
\end{verbatim}

\subsection{Invalidaci\'on de cach\'e}

Todos los hooks utilizan \texttt{useQueryClient} para invalidar la
cach\'e despu\'es de cada operaci\'on, garantizando que la interfaz
se actualice con los datos m\'as recientes:

\begin{verbatim}
onSuccess: (_, variables) => {
  queryClient.invalidateQueries({
    queryKey: ['event-registration',
      variables.userId, variables.eventId],
  });
  queryClient.invalidateQueries({
    queryKey: ['event-registration-counts'],
  });
  queryClient.invalidateQueries({
    queryKey: ['my-event-registrations',
      variables.userId],
  });
},
\end{verbatim}
